%  LaTeX support: latex@mdpi.com
%  For support, please attach all files needed for compiling as well as the log file, and specify your operating system, LaTeX version, and LaTeX editor.

%=================================================================
% pandoc conditionals added to preserve backwards compatibility with previous versions of rticles

\documentclass[notspecified,article,submit,moreauthors,pdftex]{Definitions/mdpi}


%% Some pieces required from the pandoc template
\setlist[itemize]{leftmargin=*,labelsep=5.8mm}
\setlist[enumerate]{leftmargin=*,labelsep=4.9mm}


%--------------------
% Class Options:
%--------------------
%----------
% journal
%----------
% Choose between the following MDPI journals:
% acoustics, actuators, addictions, admsci, adolescents, aerobiology, aerospace, agriculture, agriengineering, agrochemicals, agronomy, ai, air, algorithms, allergies, alloys, analytica, analytics, anatomia, animals, antibiotics, antibodies, antioxidants, applbiosci, appliedchem, appliedmath, applmech, applmicrobiol, applnano, applsci, aquacj, architecture, arm, arthropoda, arts, asc, asi, astronomy, atmosphere, atoms, audiolres, automation, axioms, bacteria, batteries, bdcc, behavsci, beverages, biochem, bioengineering, biologics, biology, biomass, biomechanics, biomed, biomedicines, biomedinformatics, biomimetics, biomolecules, biophysica, biosensors, biotech, birds, bloods, blsf, brainsci, breath, buildings, businesses, cancers, carbon, cardiogenetics, catalysts, cells, ceramics, challenges, chemengineering, chemistry, chemosensors, chemproc, children, chips, cimb, civileng, cleantechnol, climate, clinpract, clockssleep, cmd, coasts, coatings, colloids, colorants, commodities, compounds, computation, computers, condensedmatter, conservation, constrmater, cosmetics, covid, crops, cryptography, crystals, csmf, ctn, curroncol, cyber, dairy, data, ddc, dentistry, dermato, dermatopathology, designs, devices, diabetology, diagnostics, dietetics, digital, disabilities, diseases, diversity, dna, drones, dynamics, earth, ebj, ecologies, econometrics, economies, education, ejihpe, electricity, electrochem, electronicmat, electronics, encyclopedia, endocrines, energies, eng, engproc, entomology, entropy, environments, environsciproc, epidemiologia, epigenomes, est, fermentation, fibers, fintech, fire, fishes, fluids, foods, forecasting, forensicsci, forests, foundations, fractalfract, fuels, future, futureinternet, futurepharmacol, futurephys, futuretransp, galaxies, games, gases, gastroent, gastrointestdisord, gels, genealogy, genes, geographies, geohazards, geomatics, geosciences, geotechnics, geriatrics, grasses, gucdd, hazardousmatters, healthcare, hearts, hemato, hematolrep, heritage, higheredu, highthroughput, histories, horticulturae, hospitals, humanities, humans, hydrobiology, hydrogen, hydrology, hygiene, idr, ijerph, ijfs, ijgi, ijms, ijns, ijpb, ijtm, ijtpp, ime, immuno, informatics, information, infrastructures, inorganics, insects, instruments, inventions, iot, j, jal, jcdd, jcm, jcp, jcs, jcto, jdb, jeta, jfb, jfmk, jimaging, jintelligence, jlpea, jmmp, jmp, jmse, jne, jnt, jof, joitmc, jor, journalmedia, jox, jpm, jrfm, jsan, jtaer, jvd, jzbg, kidneydial, kinasesphosphatases, knowledge, land, languages, laws, life, liquids, literature, livers, logics, logistics, lubricants, lymphatics, machines, macromol, magnetism, magnetochemistry, make, marinedrugs, materials, materproc, mathematics, mca, measurements, medicina, medicines, medsci, membranes, merits, metabolites, metals, meteorology, methane, metrology, micro, microarrays, microbiolres, micromachines, microorganisms, microplastics, minerals, mining, modelling, molbank, molecules, mps, msf, mti, muscles, nanoenergyadv, nanomanufacturing,\gdef\@continuouspages{yes}} nanomaterials, ncrna, ndt, network, neuroglia, neurolint, neurosci, nitrogen, notspecified, %%nri, nursrep, nutraceuticals, nutrients, obesities, oceans, ohbm, onco, %oncopathology, optics, oral, organics, organoids, osteology, oxygen, parasites, parasitologia, particles, pathogens, pathophysiology, pediatrrep, pharmaceuticals, pharmaceutics, pharmacoepidemiology,\gdef\@ISSN{2813-0618}\gdef\@continuous pharmacy, philosophies, photochem, photonics, phycology, physchem, physics, physiologia, plants, plasma, platforms, pollutants, polymers, polysaccharides, poultry, powders, preprints, proceedings, processes, prosthesis, proteomes, psf, psych, psychiatryint, psychoactives, publications, quantumrep, quaternary, qubs, radiation, reactions, receptors, recycling, regeneration, religions, remotesensing, reports, reprodmed, resources, rheumato, risks, robotics, ruminants, safety, sci, scipharm, sclerosis, seeds, sensors, separations, sexes, signals, sinusitis, skins, smartcities, sna, societies, socsci, software, soilsystems, solar, solids, spectroscj, sports, standards, stats, std, stresses, surfaces, surgeries, suschem, sustainability, symmetry, synbio, systems, targets, taxonomy, technologies, telecom, test, textiles, thalassrep, thermo, tomography, tourismhosp, toxics, toxins, transplantology, transportation, traumacare, traumas, tropicalmed, universe, urbansci, uro, vaccines, vehicles, venereology, vetsci, vibration, virtualworlds, viruses, vision, waste, water, wem, wevj, wind, women, world, youth, zoonoticdis 
% For posting an early version of this manuscript as a preprint, you may use "preprints" as the journal. Changing "submit" to "accept" before posting will remove line numbers.

%---------
% article
%---------
% The default type of manuscript is "article", but can be replaced by: 
% abstract, addendum, article, book, bookreview, briefreport, casereport, comment, commentary, communication, conferenceproceedings, correction, conferencereport, entry, expressionofconcern, extendedabstract, datadescriptor, editorial, essay, erratum, hypothesis, interestingimage, obituary, opinion, projectreport, reply, retraction, review, perspective, protocol, shortnote, studyprotocol, systematicreview, supfile, technicalnote, viewpoint, guidelines, registeredreport, tutorial
% supfile = supplementary materials

%----------
% submit
%----------
% The class option "submit" will be changed to "accept" by the Editorial Office when the paper is accepted. This will only make changes to the frontpage (e.g., the logo of the journal will get visible), the headings, and the copyright information. Also, line numbering will be removed. Journal info and pagination for accepted papers will also be assigned by the Editorial Office.

%------------------
% moreauthors
%------------------
% If there is only one author the class option oneauthor should be used. Otherwise use the class option moreauthors.

%---------
% pdftex
%---------
% The option pdftex is for use with pdfLaTeX. Remove "pdftex" for (1) compiling with LaTeX & dvi2pdf (if eps figures are used) or for (2) compiling with XeLaTeX.

%=================================================================
% MDPI internal commands - do not modify
\firstpage{1} 
\makeatletter 
\setcounter{page}{\@firstpage} 
\makeatother
\pubvolume{1}
\issuenum{1}
\articlenumber{0}
\pubyear{2024}
\copyrightyear{2024}
%\externaleditor{Academic Editor: Firstname Lastname}
\datereceived{ } 
\daterevised{ } % Comment out if no revised date
\dateaccepted{ } 
\datepublished{ } 
%\datecorrected{} % For corrected papers: "Corrected: XXX" date in the original paper.
%\dateretracted{} % For corrected papers: "Retracted: XXX" date in the original paper.
\hreflink{https://doi.org/} % If needed use \linebreak
%\doinum{}
%\pdfoutput=1 % Uncommented for upload to arXiv.org
%\CorrStatement{yes}  % For updates


%=================================================================
% Add packages and commands here. The following packages are loaded in our class file: fontenc, inputenc, calc, indentfirst, fancyhdr, graphicx, epstopdf, lastpage, ifthen, float, amsmath, amssymb, lineno, setspace, enumitem, mathpazo, booktabs, titlesec, etoolbox, tabto, xcolor, colortbl, soul, multirow, microtype, tikz, totcount, changepage, attrib, upgreek, array, tabularx, pbox, ragged2e, tocloft, marginnote, marginfix, enotez, amsthm, natbib, hyperref, cleveref, scrextend, url, geometry, newfloat, caption, draftwatermark, seqsplit
% cleveref: load \crefname definitions after \begin{document}

%=================================================================
% Please use the following mathematics environments: Theorem, Lemma, Corollary, Proposition, Characterization, Property, Problem, Example, ExamplesandDefinitions, Hypothesis, Remark, Definition, Notation, Assumption
%% For proofs, please use the proof environment (the amsthm package is loaded by the MDPI class).

%=================================================================
% Full title of the paper (Capitalized)
\Title{Proyecto Análisis Exploratorio de Datos}

% MDPI internal command: Title for citation in the left column
\TitleCitation{Proyecto Análisis Exploratorio de Datos}

% Author Orchid ID: enter ID or remove command
%\newcommand{\orcidauthorA}{0000-0000-0000-000X} % Add \orcidA{} behind the author's name
%\newcommand{\orcidauthorB}{0000-0000-0000-000X} % Add \orcidB{} behind the author's name


% Authors, for the paper (add full first names)
\Author{María de los Ángeles Díaz Castro$^{}$, Irene Barba La
Orden$^{}$, Ana Blasco Vega$^{}$}


%\longauthorlist{yes}


% MDPI internal command: Authors, for metadata in PDF
\AuthorNames{María de los Ángeles Díaz Castro, Irene Barba La Orden, Ana
Blasco Vega}

% MDPI internal command: Authors, for citation in the left column

% Affiliations / Addresses (Add [1] after \address if there is only one affiliation.)
\address{%
}

% Contact information of the corresponding author
\corres{Correspondence: }

% Current address and/or shared authorship








% The commands \thirdnote{} till \eighthnote{} are available for further notes

% Simple summary
\simplesumm{Preparación y análisis exploratorio del tráfico diario de
carreteras en Valencia 2025, obtenido de la página oficial de la
Generalitat Valenciana.}

%\conference{} % An extended version of a conference paper

% Abstract (Do not insert blank lines, i.e. \\)
\abstract{Este proyecto tiene como objetivo analizar el tráfico en
Valencia durante 2025, utilizando datos del dataset público
proporcionado por la Generalitat Valenciana. Se consideran 18 variables,
de las que se incluyen el número de vehículos que transitan por cada
tramo, el método de captación o las coordenadas iniciales y finales de
cada registro, entre otras. En total, se han recopilado 15.235
observaciones asociadas a distintos momentos en los que se registró
tráfico. Los datos se reestructuraron para facilitar un análisis
exploratorio previo, identificando patrones de circulación y posibles
anomalías, tratando de ofrecer información valiosa para mejorar la
gestión del tráfico y la planificación de las carreteras en Valencia.}


% Keywords

% The fields PACS, MSC, and JEL may be left empty or commented out if not applicable
%\PACS{J0101}
%\MSC{}
%\JEL{}

%%%%%%%%%%%%%%%%%%%%%%%%%%%%%%%%%%%%%%%%%%
% Only for the journal Diversity
%\LSID{\url{http://}}

%%%%%%%%%%%%%%%%%%%%%%%%%%%%%%%%%%%%%%%%%%
% Only for the journal Applied Sciences

%%%%%%%%%%%%%%%%%%%%%%%%%%%%%%%%%%%%%%%%%%

%%%%%%%%%%%%%%%%%%%%%%%%%%%%%%%%%%%%%%%%%%
% Only for the journal Data



%%%%%%%%%%%%%%%%%%%%%%%%%%%%%%%%%%%%%%%%%%
% Only for the journal Toxins


%%%%%%%%%%%%%%%%%%%%%%%%%%%%%%%%%%%%%%%%%%
% Only for the journal Encyclopedia


%%%%%%%%%%%%%%%%%%%%%%%%%%%%%%%%%%%%%%%%%%
% Only for the journal Advances in Respiratory Medicine
%\addhighlights{yes}
%\renewcommand{\addhighlights}{%

%\noindent This is an obligatory section in “Advances in Respiratory Medicine”, whose goal is to increase the discoverability and readability of the article via search engines and other scholars. Highlights should not be a copy of the abstract, but a simple text allowing the reader to quickly and simplified find out what the article is about and what can be cited from it. Each of these parts should be devoted up to 2~bullet points.\vspace{3pt}\\
%\textbf{What are the main findings?}
% \begin{itemize}[labelsep=2.5mm,topsep=-3pt]
% \item First bullet.
% \item Second bullet.
% \end{itemize}\vspace{3pt}
%\textbf{What is the implication of the main finding?}
% \begin{itemize}[labelsep=2.5mm,topsep=-3pt]
% \item First bullet.
% \item Second bullet.
% \end{itemize}
%}


%%%%%%%%%%%%%%%%%%%%%%%%%%%%%%%%%%%%%%%%%%

% Pandoc syntax highlighting
\usepackage{color}
\usepackage{fancyvrb}
\newcommand{\VerbBar}{|}
\newcommand{\VERB}{\Verb[commandchars=\\\{\}]}
\DefineVerbatimEnvironment{Highlighting}{Verbatim}{commandchars=\\\{\}}
% Add ',fontsize=\small' for more characters per line
\usepackage{framed}
\definecolor{shadecolor}{RGB}{248,248,248}
\newenvironment{Shaded}{\begin{snugshade}}{\end{snugshade}}
\newcommand{\AlertTok}[1]{\textcolor[rgb]{0.94,0.16,0.16}{#1}}
\newcommand{\AnnotationTok}[1]{\textcolor[rgb]{0.56,0.35,0.01}{\textbf{\textit{#1}}}}
\newcommand{\AttributeTok}[1]{\textcolor[rgb]{0.13,0.29,0.53}{#1}}
\newcommand{\BaseNTok}[1]{\textcolor[rgb]{0.00,0.00,0.81}{#1}}
\newcommand{\BuiltInTok}[1]{#1}
\newcommand{\CharTok}[1]{\textcolor[rgb]{0.31,0.60,0.02}{#1}}
\newcommand{\CommentTok}[1]{\textcolor[rgb]{0.56,0.35,0.01}{\textit{#1}}}
\newcommand{\CommentVarTok}[1]{\textcolor[rgb]{0.56,0.35,0.01}{\textbf{\textit{#1}}}}
\newcommand{\ConstantTok}[1]{\textcolor[rgb]{0.56,0.35,0.01}{#1}}
\newcommand{\ControlFlowTok}[1]{\textcolor[rgb]{0.13,0.29,0.53}{\textbf{#1}}}
\newcommand{\DataTypeTok}[1]{\textcolor[rgb]{0.13,0.29,0.53}{#1}}
\newcommand{\DecValTok}[1]{\textcolor[rgb]{0.00,0.00,0.81}{#1}}
\newcommand{\DocumentationTok}[1]{\textcolor[rgb]{0.56,0.35,0.01}{\textbf{\textit{#1}}}}
\newcommand{\ErrorTok}[1]{\textcolor[rgb]{0.64,0.00,0.00}{\textbf{#1}}}
\newcommand{\ExtensionTok}[1]{#1}
\newcommand{\FloatTok}[1]{\textcolor[rgb]{0.00,0.00,0.81}{#1}}
\newcommand{\FunctionTok}[1]{\textcolor[rgb]{0.13,0.29,0.53}{\textbf{#1}}}
\newcommand{\ImportTok}[1]{#1}
\newcommand{\InformationTok}[1]{\textcolor[rgb]{0.56,0.35,0.01}{\textbf{\textit{#1}}}}
\newcommand{\KeywordTok}[1]{\textcolor[rgb]{0.13,0.29,0.53}{\textbf{#1}}}
\newcommand{\NormalTok}[1]{#1}
\newcommand{\OperatorTok}[1]{\textcolor[rgb]{0.81,0.36,0.00}{\textbf{#1}}}
\newcommand{\OtherTok}[1]{\textcolor[rgb]{0.56,0.35,0.01}{#1}}
\newcommand{\PreprocessorTok}[1]{\textcolor[rgb]{0.56,0.35,0.01}{\textit{#1}}}
\newcommand{\RegionMarkerTok}[1]{#1}
\newcommand{\SpecialCharTok}[1]{\textcolor[rgb]{0.81,0.36,0.00}{\textbf{#1}}}
\newcommand{\SpecialStringTok}[1]{\textcolor[rgb]{0.31,0.60,0.02}{#1}}
\newcommand{\StringTok}[1]{\textcolor[rgb]{0.31,0.60,0.02}{#1}}
\newcommand{\VariableTok}[1]{\textcolor[rgb]{0.00,0.00,0.00}{#1}}
\newcommand{\VerbatimStringTok}[1]{\textcolor[rgb]{0.31,0.60,0.02}{#1}}
\newcommand{\WarningTok}[1]{\textcolor[rgb]{0.56,0.35,0.01}{\textbf{\textit{#1}}}}

% tightlist command for lists without linebreak
\providecommand{\tightlist}{%
  \setlength{\itemsep}{0pt}\setlength{\parskip}{0pt}}


% Add imagehandling



\usepackage{longtable}

\begin{document}



%%%%%%%%%%%%%%%%%%%%%%%%%%%%%%%%%%%%%%%%%%

\section{Introducción.}\label{introducciuxf3n.}

El tráfico de vehículos va más allá de los números, muestra cómo
Valencia se mueve, organiza y enfrenta cada día sus desafíos de
desplazamiento. Este proyecto se centra en analizar la intensidad del
tráfico en las carreteras de la Comunidad Valenciana durante 2025, a
partir de datos públicos proporcionados por la Generalitat Valenciana. A
través de un análisis exploratorio de 15.235 observaciones, se busca
descubrir patrones de circulación, identificar posibles anomalías y
generar información que sea útil para la gestión del transporte. Las
visualizaciones y estadísticas que se presentan a lo largo del documento
permiten observar de manera clara cómo se mueve Valencia por sus
carreteras, ofreciendo una visión tanto cuantitativa como práctica del
tráfico en 2025.

\section{Importación de los datos.}\label{importaciuxf3n-de-los-datos.}

Para importar los datos, primero se descarga el fichero y se verifica
que los datos estén en el directorio adecuado. De esta forma, haciendo
uso de la función de importación en R, obtenemos el dataset. Es
importante especificar el encoding, ya que, de lo contrario, algunos
caracteres pueden no reconocerse correctamente.

\begin{Shaded}
\begin{Highlighting}[]
\NormalTok{filename }\OtherTok{\textless{}{-}}\StringTok{\textquotesingle{}../data/intensidad\_vehiculos\_en\_carretera\_2025.csv\textquotesingle{}}
\NormalTok{vehiculos }\OtherTok{\textless{}{-}} \FunctionTok{read.csv}\NormalTok{(filename, }\AttributeTok{sep=}\StringTok{";"}\NormalTok{,}\AttributeTok{encoding =} \StringTok{"UTF{-}8"}\NormalTok{)}
\end{Highlighting}
\end{Shaded}

Las 18 variables del dataset son las siguientes:

\begin{enumerate}
\item \textbf{TRAM}: Código del tramo al que se le asocia el tráfico.
\subitem  Primer y segundo dígito: código CV- de la carretera; tres últimos dígitos: número de tramo.
\item \textbf{DIA}: Fecha de los datos.
\item \textbf{INTA}: Total de vehículos que transitan en el sentido ascendente de la kilometración de la carretera.
\item \textbf{INTAP}: Total de vehículos pesados que transitan en el sentido ascendente de la kilometración de la carretera.
\item \textbf{INTD}: Total de vehículos que transitan en el sentido descendente de la kilometración de la carretera.
\item \textbf{INTDP}: Total de vehículos pesados que transitan en el sentido descendente de la kilometración de la carretera.
\item \textbf{INTT}: Total de vehículos que pasan por el tramo analizado (este dato será la suma de los dos sentidos anteriores).
\item \textbf{INTTP}: Total de vehículos pesados que pasan por el tramo analizado (este dato será la suma de los dos sentidos anteriores).
\item \textbf{TIPUSV}: Tipo de vía (Autovía, Convencional o Desdoblada).
\item \textbf{TIPUSE}: Procedimiento de captura de tráfico (Espiras, Gomas Neumáticas, Cámaras o Radar).
\item \textbf{PKE}: Posición exacta, en metros, donde se ha captado el tráfico, tomando como referencia la placa de kilometraje anterior en sentido ascendente.
\item \textbf{PKECOOR}: Coordenadas geodésicas en latitud y longitud del punto donde se mide el tráfico.
\item \textbf{PKINI}: Posición inicial del tramo al que se asocia el tráfico, en metros, tomando como referencia la placa kilométrica anterior en sentido ascendente (tráfico invariante en todo el tramo).
\item \textbf{REFINI}: Ubicación del inicio del tramo al que se asocia el tráfico.
\item \textbf{PKINICOOR}: Coordenadas geodésicas en latitud y longitud del punto inicial del tramo al que se asocia el tráfico (tráfico invariante en todo el tramo).
\item \textbf{PKFI}: Posición final del tramo al que se asocia el tráfico, en metros,tomando como referencia la placa kilométrica anterior en sentido ascendente (tráfico invariante en todo el tramo).
\item \textbf{REFFI}: Ubicación del final del tramo al que se asocia el tráfico.
\item \textbf{PKFICOOR}: Coordenadas geodésicas en latitud y longitud del punto final del tramo al que se asocia el tráfico (tráfico invariante en todo el tramo).
\end{enumerate}

Es importante revisar el tipo de cada variable para comprender mejor la
estructura del dataset, para ello se utilizan las funciones sapply y
class:

\begin{Shaded}
\begin{Highlighting}[]
\FunctionTok{sapply}\NormalTok{(vehiculos, class)}
\end{Highlighting}
\end{Shaded}

\begin{verbatim}
##        TRAM         DIA        INTA       INTAP        INTD       INTDP 
##   "integer" "character"   "numeric"   "numeric"   "numeric"   "numeric" 
##        INTT       INTTP      TIPUSV      TIPUSE         PKE     PKECOOR 
##   "numeric"   "numeric" "character" "character"   "integer" "character" 
##       PKINI      REFINI   PKINICOOR        PKFI       REFFI    PKFICOOR 
##   "integer" "character" "character"   "integer" "character" "character"
\end{verbatim}

Se puede observar que \(DIA\) es de tipo character, por lo que es
necesario convertirla a formato de fecha.

\begin{Shaded}
\begin{Highlighting}[]
\NormalTok{vehiculos}\SpecialCharTok{$}\NormalTok{DIA }\OtherTok{\textless{}{-}} \FunctionTok{as.Date}\NormalTok{(vehiculos}\SpecialCharTok{$}\NormalTok{DIA, }\AttributeTok{format =} \StringTok{\textquotesingle{}\%d/\%m/\%Y\textquotesingle{}}\NormalTok{)}
\FunctionTok{class}\NormalTok{(vehiculos}\SpecialCharTok{$}\NormalTok{DIA)}
\end{Highlighting}
\end{Shaded}

\begin{verbatim}
## [1] "Date"
\end{verbatim}

\section{Conversión de los datos a estructura
Tidy.}\label{conversiuxf3n-de-los-datos-a-estructura-tidy.}

La estructura tidy se caracteriza por tener cada variable en una
columna, cada observación en una fila y nombres de columnas que sean
representativos. Veamos que ocure en nuestro dataset.

La columna \(TRAM\) contiene información del código CV- de la carretera
(dos primeros dígitos) y del número de tramo(tres últimos dígitos). Por
tanto, primero verificaremos que todos los valores de esta columna
tengan cinco dígitos. Si es así, procederemos a dividirla en dos
columnas separadas, que se llamarán \(Código_CV\) y \(Num_tramo\).

\begin{Shaded}
\begin{Highlighting}[]
\FunctionTok{library}\NormalTok{(tidyr)}
\end{Highlighting}
\end{Shaded}

\begin{verbatim}
## Warning: package 'tidyr' was built under R version 4.4.2
\end{verbatim}

\begin{Shaded}
\begin{Highlighting}[]
\FunctionTok{library}\NormalTok{(dplyr)}
\end{Highlighting}
\end{Shaded}

\begin{verbatim}
## Warning: package 'dplyr' was built under R version 4.4.3
\end{verbatim}

\begin{verbatim}
## 
## Adjuntando el paquete: 'dplyr'
\end{verbatim}

\begin{verbatim}
## The following objects are masked from 'package:stats':
## 
##     filter, lag
\end{verbatim}

\begin{verbatim}
## The following objects are masked from 'package:base':
## 
##     intersect, setdiff, setequal, union
\end{verbatim}

\begin{Shaded}
\begin{Highlighting}[]
\FunctionTok{all}\NormalTok{(}\FunctionTok{length}\NormalTok{(vehiculos}\SpecialCharTok{$}\NormalTok{TRAM }\SpecialCharTok{==} \DecValTok{5}\NormalTok{))}
\end{Highlighting}
\end{Shaded}

\begin{verbatim}
## [1] TRUE
\end{verbatim}

\begin{Shaded}
\begin{Highlighting}[]
\NormalTok{vehiculos\_tidy }\OtherTok{\textless{}{-}}\NormalTok{ vehiculos }\SpecialCharTok{\%\textgreater{}\%}
  \FunctionTok{separate}\NormalTok{(}\AttributeTok{col =}\NormalTok{ TRAM, }\AttributeTok{into =} \FunctionTok{c}\NormalTok{(}\StringTok{\textquotesingle{}Codigo\_CV\textquotesingle{}}\NormalTok{, }\StringTok{\textquotesingle{}Num\_tramo\textquotesingle{}}\NormalTok{), }\AttributeTok{sep =} \DecValTok{2}\NormalTok{)}
\end{Highlighting}
\end{Shaded}

Las variables \(INTA\), \(INTAP\), \(INTD\), \(INTDP\), \(INTT\) y
\(INTTP\) contienen información sobre el tipo de vehículo, si es pesado
o no, y el sentido del tráfico, ascendente o descendente. Con el fin de
estructurar mejor estos datos, aplicamos pivot\_longer(), generando dos
nuevas variables, \(Sentido\) y \(Tipo_Vehiculo\).

En la columna \(Tipo_Vehiculo\), los valores en blanco se rellenan con
`T', que representa cualquier vehículo, mientras que `P' indica
vehículos pesados.

Por otro lado, las columnas \(INTT\) y \(INTTP\) representan la
intensidad total del tráfico, es decir, la suma de los vehículos que
circulan en ambos sentidos (ascendente y descendente), diferenciando si
son pesados o no. Dado que esta información puede obtenerse a partir de
las demás variables, eliminaremos estas columnas y, en caso de ser
necesario, podremos recalcular sus valores sumando las correspondientes
columnas de sentido.

\begin{Shaded}
\begin{Highlighting}[]
\NormalTok{vehiculos\_tidy }\OtherTok{\textless{}{-}}\NormalTok{ vehiculos\_tidy }\SpecialCharTok{\%\textgreater{}\%}
  \FunctionTok{pivot\_longer}\NormalTok{(}\AttributeTok{col =} \FunctionTok{starts\_with}\NormalTok{(}\StringTok{\textquotesingle{}INT\textquotesingle{}}\NormalTok{), }
    \AttributeTok{names\_to =} \FunctionTok{c}\NormalTok{(}\StringTok{\textquotesingle{}Sentido\textquotesingle{}}\NormalTok{, }\StringTok{\textquotesingle{}Tipo\_Vehiculo\textquotesingle{}}\NormalTok{), }
    \AttributeTok{names\_pattern =} \StringTok{\textquotesingle{}\^{}INT([ADT])\{1\}(P?)$\textquotesingle{}}\NormalTok{, }\AttributeTok{values\_to =} \StringTok{\textquotesingle{}Intensidad\textquotesingle{}}\NormalTok{)}\SpecialCharTok{\%\textgreater{}\%}
  \FunctionTok{mutate}\NormalTok{(}\AttributeTok{Tipo\_Vehiculo =} \FunctionTok{ifelse}\NormalTok{(Tipo\_Vehiculo }\SpecialCharTok{==} \StringTok{\textquotesingle{}P\textquotesingle{}}\NormalTok{, }\StringTok{\textquotesingle{}P\textquotesingle{}}\NormalTok{, }\StringTok{\textquotesingle{}T\textquotesingle{}}\NormalTok{))}\SpecialCharTok{\%\textgreater{}\%}
  \FunctionTok{filter}\NormalTok{(Sentido }\SpecialCharTok{!=} \StringTok{\textquotesingle{}T\textquotesingle{}}\NormalTok{)}
\end{Highlighting}
\end{Shaded}

Del mismo modo, las columnas \(PKECOOR\), \(PKINICOOR\) y \(PKFICOOR\)
contienen las coordenadas geográficas, diferenciando entre la posición
inicial, final y exacta. Para reorganizar esta información, aplicamos
pivot\_longer(), que genera dos nuevas columnas: \(Coordenadas\), donde
se indica el tipo de coordenada (inicial, final o exacta), y
\(Valor_coordenada\), que almacena su valor correspondiente.

Además, observamos que cada coordenada incluye la latitud y la longitud
separadas por un punto y coma. Por ello, utilizamos la función
separate() para dividir \(Valor_coordenada\) en dos columnas
adicionales: \(Latitud\) y \(Longitud\).

\begin{Shaded}
\begin{Highlighting}[]
\NormalTok{vehiculos\_tidy }\OtherTok{\textless{}{-}}\NormalTok{ vehiculos\_tidy }\SpecialCharTok{\%\textgreater{}\%}
  \FunctionTok{pivot\_longer}\NormalTok{(}\AttributeTok{col =} \FunctionTok{ends\_with}\NormalTok{(}\StringTok{\textquotesingle{}COOR\textquotesingle{}}\NormalTok{), }\AttributeTok{names\_to =} \StringTok{\textquotesingle{}Coordenadas\textquotesingle{}}\NormalTok{, }
               \AttributeTok{names\_pattern =} \StringTok{\textquotesingle{}\^{}PK(E|INI|FI)COOR$\textquotesingle{}}\NormalTok{,  }
               \AttributeTok{values\_to =} \StringTok{\textquotesingle{}Valor\_coordenada\textquotesingle{}}\NormalTok{)}\SpecialCharTok{\%\textgreater{}\%}
  \FunctionTok{separate}\NormalTok{(}\AttributeTok{col =}\NormalTok{ Valor\_coordenada, }\AttributeTok{into =} \FunctionTok{c}\NormalTok{(}\StringTok{\textquotesingle{}Latitud\textquotesingle{}}\NormalTok{, }\StringTok{\textquotesingle{}Longitud\textquotesingle{}}\NormalTok{), }\AttributeTok{sep=}\StringTok{\textquotesingle{};\textquotesingle{}}\NormalTok{)}\SpecialCharTok{\%\textgreater{}\%}
  \FunctionTok{pivot\_longer}\NormalTok{(}\AttributeTok{col =} \FunctionTok{c}\NormalTok{(Latitud, Longitud), }\AttributeTok{names\_to =} \StringTok{\textquotesingle{}Tipo\_Coordenada\textquotesingle{}}\NormalTok{, }
               \AttributeTok{values\_to =} \StringTok{\textquotesingle{}Valor\_Coordenada\textquotesingle{}}\NormalTok{) }
\end{Highlighting}
\end{Shaded}

En cuanto a las \(REFINI\) y \(REFFI\) recogen las ubicaciones inicial y
final del tramo donde se ha detectado el tráfico. Para reorganizar esta
información, aplicamos pivot\_longer(), generando dos nuevas columnas:
\(Tipo_Ubicacion\), que especifica si se trata de la ubicación inicial o
final, y \(Ubicacion\), que que almacena el nombre de la vía o
localización.

\begin{Shaded}
\begin{Highlighting}[]
\NormalTok{vehiculos\_tidy }\OtherTok{\textless{}{-}}\NormalTok{ vehiculos\_tidy }\SpecialCharTok{\%\textgreater{}\%}
  \FunctionTok{pivot\_longer}\NormalTok{(}\AttributeTok{col =} \FunctionTok{starts\_with}\NormalTok{(}\StringTok{\textquotesingle{}REF\textquotesingle{}}\NormalTok{), }\AttributeTok{names\_pattern =} \StringTok{\textquotesingle{}\^{}REF(INI|FI)\textquotesingle{}}\NormalTok{,}
               \AttributeTok{names\_to =} \StringTok{\textquotesingle{}Tipo\_Ubicacion\textquotesingle{}}\NormalTok{, }\AttributeTok{values\_to =} \StringTok{\textquotesingle{}Ubicacion\textquotesingle{}}\NormalTok{)}
\end{Highlighting}
\end{Shaded}

Por último, las variables \(PKINI\), \(PKFI\) y \(PKE\) indican la
posición inicial, final y exacta donde se registró el tráfico, tomando
como referencia la placa kilométrica anterior en sentido ascendente,
suponiendo que el tráfico es invariante en todo el tramo.

Mediante la función pivot\_longer(), se crean dos nuevas columnas:
\(Tipo_Posicion_Punto_Quilometrico\), que especifica si la posición
corresponde al punto inicial o final, y \(Valor_Punto_Quilometrico\),
que contiene el valor asociado a dicha posición.

\begin{Shaded}
\begin{Highlighting}[]
\NormalTok{vehiculos\_tidy }\OtherTok{\textless{}{-}}\NormalTok{ vehiculos\_tidy }\SpecialCharTok{\%\textgreater{}\%}
   \FunctionTok{pivot\_longer}\NormalTok{(}\AttributeTok{col =} \FunctionTok{matches}\NormalTok{(}\StringTok{\textquotesingle{}\^{}PK(INI|FI|E)\textquotesingle{}}\NormalTok{), }\AttributeTok{names\_pattern =} \StringTok{\textquotesingle{}\^{}PK(INI|FI|E)\textquotesingle{}}\NormalTok{,}
                \AttributeTok{names\_to =} \StringTok{\textquotesingle{}Tipo\_Posicion\_Punto\_Quilometrico\textquotesingle{}}\NormalTok{, }
                \AttributeTok{values\_to =} \StringTok{\textquotesingle{}Valor\_Punto\_Quilometrico\textquotesingle{}}\NormalTok{)}
\end{Highlighting}
\end{Shaded}

La columna \(Valor_Coordenada\) contiene los valores de latitud y
longitud, donde se utilizan comas como separador decimal. Por este
motivo, R detecta la columna como tipo character, cuando en realidad
debería ser numérica. Para solucionarlo, usamos la función sub() para
reemplazar las comas por puntos, lo que nos permite convertir
correctamente los valores a tipo numérico. Veámoslo:

\begin{Shaded}
\begin{Highlighting}[]
\FunctionTok{str}\NormalTok{(vehiculos\_tidy}\SpecialCharTok{$}\NormalTok{Valor\_Coordenada)}
\end{Highlighting}
\end{Shaded}

\begin{verbatim}
##  chr [1:2193840] "39,88828171" "39,88828171" "39,88828171" "39,88828171" ...
\end{verbatim}

\begin{Shaded}
\begin{Highlighting}[]
\NormalTok{vehiculos\_tidy}\SpecialCharTok{$}\NormalTok{Valor\_Coordenada }\OtherTok{\textless{}{-}} \FunctionTok{sub}\NormalTok{(}\StringTok{\textquotesingle{},\textquotesingle{}}\NormalTok{, }\StringTok{\textquotesingle{}.\textquotesingle{}}\NormalTok{, vehiculos\_tidy}\SpecialCharTok{$}\NormalTok{Valor\_Coordenada)}
\NormalTok{vehiculos\_tidy}\SpecialCharTok{$}\NormalTok{Valor\_Coordenada }\OtherTok{\textless{}{-}}\FunctionTok{as.numeric}\NormalTok{(vehiculos\_tidy}\SpecialCharTok{$}\NormalTok{Valor\_Coordenada)}
\FunctionTok{str}\NormalTok{(vehiculos\_tidy}\SpecialCharTok{$}\NormalTok{Valor\_Coordenada)}
\end{Highlighting}
\end{Shaded}

\begin{verbatim}
##  num [1:2193840] 39.9 39.9 39.9 39.9 39.9 ...
\end{verbatim}

Solo falta asegurarnos de que los nombres de las columnas sean
representativos; para ello, utilizamos la función rename() y
reemplazamos algunos valores abreviados por descripciones completas
mediante la función ifelse(). Por ejemplo, en las variables relacionadas
con coordenadas y ubicación, los códigos INI, E y FI se sustituyeron por
Inicial, Exacta y Final; en Sentido, los valores A y D pasaron a
Ascendente y Descendente; y en Tipo\_Vehiculo, T y P se cambiaron por
Todos y Pesados. Veámoslo:

\begin{Shaded}
\begin{Highlighting}[]
\NormalTok{vehiculos\_tidy }\OtherTok{\textless{}{-}}\NormalTok{ vehiculos\_tidy }\SpecialCharTok{\%\textgreater{}\%}
\NormalTok{  dplyr}\SpecialCharTok{::}\FunctionTok{rename}\NormalTok{(}\AttributeTok{Fecha =}\NormalTok{ DIA, }\AttributeTok{Tipo\_Via =}\NormalTok{ TIPUSV, }\AttributeTok{Metodologia\_Captacion=}\NormalTok{ TIPUSE)}

\CommentTok{\# Rescribimos algunos valores:}
\NormalTok{vehiculos\_tidy}\SpecialCharTok{$}\NormalTok{Sentido }\OtherTok{\textless{}{-}}
\NormalTok{  vehiculos\_tidy}\SpecialCharTok{$}\NormalTok{Sentido }\SpecialCharTok{\%\textgreater{}\%}
  \FunctionTok{replace}\NormalTok{(. }\SpecialCharTok{==} \StringTok{"A"}\NormalTok{, }\StringTok{"Ascendente"}\NormalTok{) }\SpecialCharTok{\%\textgreater{}\%}
  \FunctionTok{replace}\NormalTok{(. }\SpecialCharTok{==} \StringTok{"D"}\NormalTok{, }\StringTok{"Descendente"}\NormalTok{)}

\NormalTok{vehiculos\_tidy}\SpecialCharTok{$}\NormalTok{Tipo\_Vehiculo }\OtherTok{\textless{}{-}}
\NormalTok{  vehiculos\_tidy}\SpecialCharTok{$}\NormalTok{Tipo\_Vehiculo }\SpecialCharTok{\%\textgreater{}\%}
  \FunctionTok{replace}\NormalTok{(. }\SpecialCharTok{==} \StringTok{"T"}\NormalTok{, }\StringTok{"Todos"}\NormalTok{) }\SpecialCharTok{\%\textgreater{}\%}
  \FunctionTok{replace}\NormalTok{(. }\SpecialCharTok{==} \StringTok{"P"}\NormalTok{, }\StringTok{"Pesados"}\NormalTok{)}

\NormalTok{vehiculos\_tidy}\SpecialCharTok{$}\NormalTok{Coordenadas }\OtherTok{\textless{}{-}}
\NormalTok{  vehiculos\_tidy}\SpecialCharTok{$}\NormalTok{Coordenadas }\SpecialCharTok{\%\textgreater{}\%}
  \FunctionTok{replace}\NormalTok{(. }\SpecialCharTok{==} \StringTok{"INI"}\NormalTok{, }\StringTok{"Inicial"}\NormalTok{) }\SpecialCharTok{\%\textgreater{}\%}
  \FunctionTok{replace}\NormalTok{(. }\SpecialCharTok{==} \StringTok{"E"}\NormalTok{, }\StringTok{"Exacta"}\NormalTok{) }\SpecialCharTok{\%\textgreater{}\%}
  \FunctionTok{replace}\NormalTok{(. }\SpecialCharTok{==} \StringTok{"FI"}\NormalTok{, }\StringTok{"Final"}\NormalTok{)}

\NormalTok{vehiculos\_tidy}\SpecialCharTok{$}\NormalTok{Tipo\_Ubicacion }\OtherTok{\textless{}{-}}
\NormalTok{  vehiculos\_tidy}\SpecialCharTok{$}\NormalTok{Tipo\_Ubicacion }\SpecialCharTok{\%\textgreater{}\%}
  \FunctionTok{replace}\NormalTok{(. }\SpecialCharTok{==} \StringTok{"INI"}\NormalTok{, }\StringTok{"Inicial"}\NormalTok{) }\SpecialCharTok{\%\textgreater{}\%}
  \FunctionTok{replace}\NormalTok{(. }\SpecialCharTok{==} \StringTok{"E"}\NormalTok{, }\StringTok{"Exacta"}\NormalTok{) }\SpecialCharTok{\%\textgreater{}\%}
  \FunctionTok{replace}\NormalTok{(. }\SpecialCharTok{==} \StringTok{"FI"}\NormalTok{, }\StringTok{"Final"}\NormalTok{)}

\NormalTok{vehiculos\_tidy}\SpecialCharTok{$}\NormalTok{Tipo\_Posicion\_Punto\_Quilometrico }\OtherTok{\textless{}{-}}
\NormalTok{  vehiculos\_tidy}\SpecialCharTok{$}\NormalTok{Tipo\_Posicion\_Punto\_Quilometrico }\SpecialCharTok{\%\textgreater{}\%}
  \FunctionTok{replace}\NormalTok{(. }\SpecialCharTok{==} \StringTok{"INI"}\NormalTok{, }\StringTok{"Inicial"}\NormalTok{) }\SpecialCharTok{\%\textgreater{}\%}
  \FunctionTok{replace}\NormalTok{(. }\SpecialCharTok{==} \StringTok{"E"}\NormalTok{, }\StringTok{"Exacta"}\NormalTok{) }\SpecialCharTok{\%\textgreater{}\%}
  \FunctionTok{replace}\NormalTok{(. }\SpecialCharTok{==} \StringTok{"FI"}\NormalTok{, }\StringTok{"Final"}\NormalTok{)}

\NormalTok{vehiculos\_tidy}\SpecialCharTok{$}\NormalTok{Tipo\_Via }\OtherTok{\textless{}{-}}
\NormalTok{  vehiculos\_tidy}\SpecialCharTok{$}\NormalTok{Tipo\_Via }\SpecialCharTok{\%\textgreater{}\%}
  \FunctionTok{replace}\NormalTok{(. }\SpecialCharTok{==} \StringTok{"Aut."}\NormalTok{, }\StringTok{"Autovia"}\NormalTok{) }\SpecialCharTok{\%\textgreater{}\%}
  \FunctionTok{replace}\NormalTok{(. }\SpecialCharTok{==} \StringTok{"Conv."}\NormalTok{, }\StringTok{"Convencional"}\NormalTok{) }\SpecialCharTok{\%\textgreater{}\%}
  \FunctionTok{replace}\NormalTok{(. }\SpecialCharTok{==} \StringTok{"Desd."}\NormalTok{, }\StringTok{"Desdoblada"}\NormalTok{)}
\end{Highlighting}
\end{Shaded}

A continuación, se muestra un fragmento de las primeras filas del
dataset, usando la función head():

\begin{Shaded}
\begin{Highlighting}[]
\FunctionTok{head}\NormalTok{(vehiculos\_tidy)}
\end{Highlighting}
\end{Shaded}

\begin{verbatim}
## # A tibble: 6 x 15
##   Codigo_CV Num_tramo Fecha      Tipo_Via Metodologia_Captacion Sentido   
##   <chr>     <chr>     <date>     <chr>    <chr>                 <chr>     
## 1 10        007       2025-03-07 Autovia  Espiras               Ascendente
## 2 10        007       2025-03-07 Autovia  Espiras               Ascendente
## 3 10        007       2025-03-07 Autovia  Espiras               Ascendente
## 4 10        007       2025-03-07 Autovia  Espiras               Ascendente
## 5 10        007       2025-03-07 Autovia  Espiras               Ascendente
## 6 10        007       2025-03-07 Autovia  Espiras               Ascendente
## # i 9 more variables: Tipo_Vehiculo <chr>, Intensidad <dbl>, Coordenadas <chr>,
## #   Tipo_Coordenada <chr>, Valor_Coordenada <dbl>, Tipo_Ubicacion <chr>,
## #   Ubicacion <chr>, Tipo_Posicion_Punto_Quilometrico <chr>,
## #   Valor_Punto_Quilometrico <int>
\end{verbatim}

%%%%%%%%%%%%%%%%%%%%%%%%%%%%%%%%%%%%%%%%%%

\vspace{6pt}

%%%%%%%%%%%%%%%%%%%%%%%%%%%%%%%%%%%%%%%%%%
%% optional
\supplementary{The following supporting information can be downloaded
at:\\
\linksupplementary{s1}, Figure S1: title; Table S1: title; Video S1:
title.}

% Only for the journal Methods and Protocols:
% If you wish to submit a video article, please do so with any other supplementary material.
% \supplementary{The following supporting information can be downloaded at: \linksupplementary{s1}, Figure S1: title; Table S1: title; Video S1: title. A supporting video article is available at doi: link.}

% Only for journal Hardware:
% If you wish to submit a video article, please do so with any other supplementary material.
% \supplementary{The following supporting information can be downloaded at: \linksupplementary{s1}, Figure S1: title; Table S1: title; Video S1: title.\vspace{6pt}\\
%\begin{tabularx}{\textwidth}{lll}
%\toprule
%\textbf{Name} & \textbf{Type} & \textbf{Description} \\
%\midrule
%S1 & Python script (.py) & Script of python source code used in XX \\
%S2 & Text (.txt) & Script of modelling code used to make Figure X \\
%S3 & Text (.txt) & Raw data from experiment X \\
%S4 & Video (.mp4) & Video demonstrating the hardware in use \\
%... & ... & ... \\
%\bottomrule
%\end{tabularx}
%}

%%%%%%%%%%%%%%%%%%%%%%%%%%%%%%%%%%%%%%%%%%




\dataavailability{Los datos usados en este estudio son públicos
disponibles en
{[}\url{https://dadesobertes.gva.es/dataset/20b91b5a-c26d-4624-bcf5-b76a74b3ec88/resource/abe57190-d53d-4d41-9266-d7c2254621ae/download/intensidad_vehiculos_en_carretera_2025.csv}{]}}

% Only for journal Nursing Reports
%\publicinvolvement{Please describe how the public (patients, consumers, carers) were involved in the research. Consider reporting against the GRIPP2 (Guidance for Reporting Involvement of Patients and the Public) checklist. If the public were not involved in any aspect of the research add: ``No public involvement in any aspect of this research''.}

% Only for journal Nursing Reports
%\guidelinesstandards{Please add a statement indicating which reporting guideline was used when drafting the report. For example, ``This manuscript was drafted against the XXX (the full name of reporting guidelines and citation) for XXX (type of research) research''. A complete list of reporting guidelines can be accessed via the equator network: \url{https://www.equator-network.org/}.}

% Only for journal Nursing Reports
%\guidelinesstandards{Please add a statement indicating which reporting guideline was used when drafting the report. For example, ``This manuscript was drafted against the XXX (the full name of reporting guidelines and citation) for XXX (type of research) research''. A complete list of reporting guidelines can be accessed via the equator network: \url{https://www.equator-network.org/}.}



%%%%%%%%%%%%%%%%%%%%%%%%%%%%%%%%%%%%%%%%%%
%% Optional

%% Only for journal Encyclopedia
%\entrylink{The Link to this entry published on the encyclopedia platform.}

\abbreviations{Abbreviations}{
The following abbreviations are used in this manuscript:\\

\noindent
\begin{tabular}{@{}ll}
MDPI & Multidisciplinary Digital Publishing Institute \\
DOAJ & Directory of open access journals \\
TLA & Three letter acronym \\
LD & linear dichroism \\
\end{tabular}
}

%%%%%%%%%%%%%%%%%%%%%%%%%%%%%%%%%%%%%%%%%%
%% Optional
\input{"appendix.tex"}
%%%%%%%%%%%%%%%%%%%%%%%%%%%%%%%%%%%%%%%%%%
\begin{adjustwidth}{-\extralength}{0cm}

%\printendnotes[custom] % Un-comment to print a list of endnotes


\reftitle{References}
\bibliography{mybibfile.bib}

% If authors have biography, please use the format below
%\section*{Short Biography of Authors}
%\bio
%{\raisebox{-0.35cm}{\includegraphics[width=3.5cm,height=5.3cm,clip,keepaspectratio]{Definitions/author1.pdf}}}
%{\textbf{Firstname Lastname} Biography of first author}
%
%\bio
%{\raisebox{-0.35cm}{\includegraphics[width=3.5cm,height=5.3cm,clip,keepaspectratio]{Definitions/author2.jpg}}}
%{\textbf{Firstname Lastname} Biography of second author}

%%%%%%%%%%%%%%%%%%%%%%%%%%%%%%%%%%%%%%%%%%
%% for journal Sci
%\reviewreports{\\
%Reviewer 1 comments and authors’ response\\
%Reviewer 2 comments and authors’ response\\
%Reviewer 3 comments and authors’ response
%}
%%%%%%%%%%%%%%%%%%%%%%%%%%%%%%%%%%%%%%%%%%
\PublishersNote{}
\end{adjustwidth}


\end{document}
